
\documentclass{article}
\usepackage{colortbl}
\usepackage{makecell}
\usepackage{multirow}
\usepackage{supertabular}

\begin{document}

\newcounter{utterance}

\centering \large Interaction Transcript for game `cladder', experiment `full\_v1.5\_default', episode 3483 with Qwen3.6{-}27B.
\vspace{24pt}

{ \footnotesize  \setcounter{utterance}{1}
\setlength{\tabcolsep}{0pt}
\begin{supertabular}{c@{$\;$}|p{.15\linewidth}@{}p{.15\linewidth}p{.15\linewidth}p{.15\linewidth}p{.15\linewidth}p{.15\linewidth}}
    \# & \multicolumn{2}{c}{Player} && \multicolumn{2}{c}{Game Master} \\
    \hline

    \theutterance \stepcounter{utterance}  
    & & & \multicolumn{4}{p{0.6\linewidth}}{
        \cellcolor[rgb]{0.9,0.9,0.9}{
            \makecell[{{p{\linewidth}}}]{
                \texttt{\tiny{[P1$\langle$GM]}}
                \texttt{You are an expert in causal inference. The following question is not a typical commonsense query, but rather a meticulously designed question created by a professor specializing in causal inference, intended to assess the students' mastery of the course content.} \\
\texttt{Imagine a self{-}contained, hypothetical world with only the following conditions, and without any unmentioned factors or causal relationships: Tanning salon treatment has a direct effect on foot size. Going to the beach has a direct effect on foot size. The overall probability of tanning salon treatment is 10\%. For people not using tanning salon treatments, the probability of large feet is 13\%. For people who went to tanning salons, the probability of large feet is 70\%. Is large feet more likely than small feet overall?} \\
\texttt{Start your answer with 'Yes' or 'No', followed by additional reasoning or evidence to support your explanation.} \\
            }
        }
    }
     \\ \\

    \theutterance \stepcounter{utterance}  
    & \multicolumn{4}{p{0.6\linewidth}}{
        \cellcolor[rgb]{0.561,0.737,0.561}{
            \makecell[{{p{\linewidth}}}]{
                \texttt{\tiny{[P1$\rangle$GM]}}
                \texttt{No} \\
\\ 
\texttt{To determine whether large feet are more likely than small feet overall, we must calculate the total probability of having large feet, denoted as \$P(L)\$. We are given the following information:} \\
\\ 
\texttt{1.  The probability of tanning salon treatment, \$P(T) = 0.10\$.} \\
\texttt{2.  Consequently, the probability of not using tanning salon treatment, \$P(\textbackslash{}neg T) = 1 {-} 0.10 = 0.90\$.} \\
\texttt{3.  The probability of large feet given no tanning salon treatment, \$P(L | \textbackslash{}neg T) = 0.13\$.} \\
\texttt{4.  The probability of large feet given tanning salon treatment, \$P(L | T) = 0.70\$.} \\
\\ 
\texttt{The problem states that "Going to the beach has a direct effect on foot size," but it does not provide any probabilities or conditional distributions regarding beach visits. However, the conditional probabilities \$P(L | T)\$ and \$P(L | \textbackslash{}neg T)\$ are given as fixed values. In causal inference problems of this structure, when conditional probabilities of the outcome given the treatment are explicitly provided, they represent the marginal probabilities of the outcome within those treatment groups, integrating out any other variables (like beach visits) that might be confounders or mediators, unless specified otherwise. Since no data is provided to adjust for beach visits, we must assume the given conditional probabilities \$P(L | T} \\
            }
        }
    }
    & & \\ \\

    \theutterance \stepcounter{utterance}  
    & & & \multicolumn{2}{p{0.3\linewidth}}{
        \cellcolor[rgb]{0.9,0.9,0.9}{
            \makecell[{{p{\linewidth}}}]{
                \texttt{\tiny{[GM$|$GM]}}
                \texttt{No} \\
            }
        }
    }
    & & \\ \\

    \theutterance \stepcounter{utterance}  
    & & & \multicolumn{2}{p{0.3\linewidth}}{
        \cellcolor[rgb]{0.9,0.9,0.9}{
            \makecell[{{p{\linewidth}}}]{
                \texttt{\tiny{[GM$|$GM]}}
                \texttt{no} \\
            }
        }
    }
    & & \\ \\

    \theutterance \stepcounter{utterance}  
    & & & \multicolumn{2}{p{0.3\linewidth}}{
        \cellcolor[rgb]{0.9,0.9,0.9}{
            \makecell[{{p{\linewidth}}}]{
                \texttt{\tiny{[GM$|$GM]}}
                \texttt{game\_result = WIN} \\
            }
        }
    }
    & & \\ \\

\end{supertabular}
}

\end{document}
